\section{Introduction}
\label{sec:introduction96}
A quantum measurement can produce only a classical label while leaving
information in a reference prepared by the experimenter.  Across two
uses, retaining this receiver, sending its classical measurement result
to a fresh source, and coherently feeding it into a later acquisition
are different operations.  This article studies their exact decision
values in a fixed operational interface.

An ordered measurement is a tuple $E=(E_y)_y$ of positive matrices with
$\sum_yE_y=I$.  Its unknown interface is
$\mathcal M_E(X)=\sum_y\tr(E_yX)|y\rangle\langle y|$;
no residual device output is accessible.  A reset cut requires the
\emph{entire} new input--reference system to be conditionally independent
of the old receiver.  The old and fresh references may be measured
jointly only after the second call.  All controls are common to the
hypotheses, and all formal histories carry a defined rule.
We use unhalved trace norms; a decision-probability perturbation carries
the corresponding factor $1/2$ for a trace-zero difference.

\subsection{General experiments and the value of a message}
For arbitrary finite priors and rewards, the branch normal form is
$\sum_h A_{yh}=\rho$, with a fresh density $b_{yh}$ on each branch.
The trace of $A_{yh}$ is a Gram-decomposition weight, not a physical
history probability.  Optimizing the fresh state and the final POVM
gives a continuation function $g_y$; receiver messages replace it by
its upper concave envelope $c_y$.  Theorems~\ref{thm:resetvariational91}
and~\ref{thm:resetdual91} give attained primal and dual values, with
finite instruments and universal Hermitian majorants.  Their
rank-resolved forms specify which complete quantum registers are
counted at each cut.

Lemma~\ref{lem:postponement96} proves the no-message reduction directly
in the Heisenberg picture.  Theorem~\ref{thm:messagecriterion96} makes
its uniform consequence exact: with the full old Schmidt support
retained, receiver messages have no value at every prescribed pure
first acquisition if and only if every fresh-rank-constrained
continuation is concave.  Failure is witnessed by two outcomes after
one first device label at some initial acquisition.  This assertion
is not a two-message bound at every preassigned initial state.
Nor may postponement bypass an old-dimension bound smaller than the
system it retains.

\subsection{A solved finite family}
We next use a weighted exact second-moment family of ordered bases.
An alternative basis is selected once and reused at both calls;
independent redraws give a different experiment with no second-moment
signal.  Theorem~\ref{thm:operationalhierarchy90} evaluates the three
unrestricted endpoint classes, and Theorem~\ref{thm:rankspectrum92}
resolves both retained-register dimensions.

For the specified biased prior and a prescribed pure initial Gram
$\rho$, Theorem~\ref{thm:spectralvalue94} gives
\[
 P_{k,\ell}^{\rm b}(\rho)
  =\frac{d+1}{2(d+1)-t^2}
       +\frac{t^2}{2[2(d+1)-t^2]}
             \chi\bigl(q^{(\min(k,\ell))}(\lambda(\rho))\bigr).
\]
The boundary-safe definitions of $\chi$ and the unique least rank-capped
majorant $q^{(r)}$, including ties, are given before that theorem in
Section~\ref{sec:spectralvalue94}.  A finite instrument on the fixed
initial receiver attains the formula.  No-message preparation instead
uses the leading truncated spectrum.  Thus at $t>0$ the message has
strict value exactly for $2\le r<\operatorname{rank}\rho$.
Theorem~\ref{thm:initialspectra93} and the subsequent rigidity estimates
describe equality and quantitative loss.  Proposition~\ref{prop:qutritwitness96}
works out a complete three-dimensional example, with exact algebraic
values and rational secular intervals.

The prior matters.  For equal priors and old receiver dimension two,
Theorem~\ref{thm:equalinitial95} gives a different initial-spectrum
formula and different optimal fresh Schmidt weights.  Its qubit
specialization covers every pure first acquisition without a receiver
message.  The full proof includes the optimizing fresh weights, the singular
boundary and finite attainment on the prescribed receiver.

\subsection{Relation to local geometry}
The same ordered-measurement interface also supports local fixed-tube
results in the linked supplement.  Those results classify support and
covariance directions, preserve the finite tube remainders, and give
allocation orders involving $N=\sum_jn_j$, $V=\sum_jn_j^2$, and hard
pathwise $Q$.  They are not used to prove the exact finite-game theorems.
Conversely the spectral formulas do not make their fixed-object
constants uniform over changing support strata.  The local theory is
retained as mathematics, not discarded; it is separated from the
finite-experiment proof route.  The independent structural article and
the wider analytic Theta programme supply no hidden hypotheses here.
